\documentclass[12pt]{article}

\usepackage[dvips,letterpaper,margin=0.75in,bottom=0.5in]{geometry}
\usepackage{cite}
\usepackage{slashed}
\usepackage{graphicx}
\usepackage{amsmath}
\usepackage{amssymb}
\usepackage{braket}
\usepackage{pythonhighlight}
\usepackage{mathtools}
\begin{document}

\newcommand{\ihbar}{\ensuremath{i \hbar}}
\newcommand{\Pss}{\ensuremath{\Psi^*}}
\newcommand{\dPsidt}{\ensuremath{ \frac{\partial \Psi}{\partial t} }}
\newcommand{\dPsidx}{\ensuremath{ \frac{\partial \Psi}{\partial x} }}
\newcommand{\ddPsidx}{\ensuremath{ \frac{\partial^2 \Psi}{\partial x^2} }}
\newcommand{\dPssdt}{\ensuremath{ \frac{\partial \Psi^*}{\partial t} }}
\newcommand{\dPssdx}{\ensuremath{ \frac{\partial \Psi^*}{\partial x} }}
\newcommand{\ddPssdx}{\ensuremath{ \frac{\partial^2 \Psi^*}{\partial x^2} }}

\newcommand{\dphidt}{\ensuremath{ \frac{d \phi}{dt} }}
\newcommand{\dpsidx}{\ensuremath{ \frac{d \psi}{dx} }}
\newcommand{\ddpsidx}{\ensuremath{ \frac{d^2 \psi}{dx^2} }}


\title{PHY 115L \\ Woods-Saxon Model for the Nucleus \\ using Variational Methods }
\author{Michael Mulhearn}

\maketitle

\section{TISE and System of Units}

Noting a good reference:  https://arxiv.org/pdf/0709.3525

\noindent
We'll adapt a system of units appropriate for modeling the atomic nucleus
with the Time-Independent
Schr\"odinger Equation (TISE):
$$\hat{H} \; \psi_n(x) = E_n \; \psi_n(x)$$
where the Hamiltonian operator is defined by:
\begin{equation}
\hat{H} = -\frac{\hbar^2}{2 M} \; \frac{d^2}{dx^2} \, + \, V(x)
\end{equation}
We set the characteristic length scale $L$ to $1~\rm fm$.
We will be considering {\bf non-periodic} wave function
(e.g. solutions to the harmonic oscillator).  To use the Fourier
series approach in this case, we take the period $a$ to be large
compared to $L$, and use the fact that the wave function approaches
zero to approximate a periodic function.

Based on our characteristic length scale $L$, a characteristic energy is:
\begin{equation}
\epsilon \equiv \frac{\hbar^2}{2 M L^2}
\end{equation}
where $M = m_p$ is the mass of the proton.  When doing nuclear and particle physics, it's helpful to use energy units, with:
\begin{eqnarray*}
\hbar c   &=& 197.3~\rm MeV fm \\
m_{\rm p} &=& 938.272~\rm MeV/c^2\\
m_{\rm n} &=& 939.565~\rm MeV/c^2\\
\end{eqnarray*}
So then:
\begin{equation}
\epsilon \equiv \frac{\hbar^2 c^2}{2 Mc^2 L^2} = 41.49~\rm MeV
\end{equation}
With these units, we have:
\begin{equation}
\frac{\hat{H}}{\epsilon} = -L^2\frac{d^2}{dx^2} + \frac{V(x)}{\epsilon}
\end{equation}

\section{Woods-Saxon Potential}

We'll use the Woods-Saxon potential:
\begin{equation}
V(r) = -\frac{V_0}{1 + \exp\left( \frac{r-R}{a} \right)}
\end{equation}
where:
\begin{equation}
R = r_0 A^{1/3} 
\end{equation}
and $A$ is the mass number.  Typical values for the parameters are:
\begin{eqnarray*}
V_0 &=& 50 ~ {\rm MeV} \\
r_0 &=& 1.25 ~ {\rm fm}\\
a   &=& 0.5 ~ {\rm fm}\\
\end{eqnarray*}
In our system of units, this becomes:
\begin{equation}
\frac{V(r/L)}{\epsilon} = -\frac{V_0/\epsilon}{1 + \exp\left( \frac{r/L-R/L}{a/L} \right)}
\end{equation}
where
$$R/L = (r_0/L) A^{1/3} $$
and
\begin{eqnarray*}
V_0/\epsilon &=& 1.205 \\
r_0/L &=& 1.25 \\
a/L   &=& 0.5 \\
\end{eqnarray*}


\section{Coulomb Potential}
When dealing with the electric charge in particle physics, it is useful to connect it to energy units via the fine structure constant:
$$\alpha = \frac{e^2}{4\pi\epsilon_0 \hbar c} = \frac{1}{137.035}$$
and
$$V(r) = Z\frac{e^2}{4\pi\epsilon_0} \frac{F(r)}{L}$$


$$F(r) \equiv  \begin{cases}
\frac{3R^2-r^2}{2R^3} & r \leq R \\
  1/r & r > R \\
\end{cases}
$$


\section{The Variational Method for Higher Energy States}

In the previous homework we found the ground state via:
$$\braket{\psi|\hat{H}|\psi} \geq E_0$$
for all $\psi$.  With knowledge that the ground state was even and the next excited state is odd, we can also determined the first excited state via:
$$\braket{\psi|\hat{H}|\psi} \geq E_1$$
for odd $\psi$.

We can generalize this to the following.  If we know that:
$$\braket{\Psi|m} = 0$$
for {\bf all} of the following:
$$m=0,1,\ldots,n$$
then we also have:  
$$\braket{\psi|\hat{H}|\psi} \geq E_{n+1}$$

We can, in principle at least, always arrange that our trial wave function is orthogonal to the previous discovered lower energy wave functions.  We use the projection operator:
$$P_m = \ket{m}\bra{m}$$
to remove any component of the trial wave function that that is not orthogonal to the (previously discovered) ground state $\ket{\Psi_0}$ to obatin:
$$\ket{\Psi} \; \longrightarrow \; (1-P_m) \ket{\Psi} =
\ket{\Psi} - \braket{0|\Psi} \; \ket{0}$$
The inner product $\braket{0|\Psi}$ is easy to calculate in momentum space.  Suppose we have discovered $d_m$ so that:
$$\ket{0} = \sum_m d_m |p_m>$$
and have:
$$\ket{\Psi} = \sum_m c_m |p_m>$$
Then:
$$\braket{0|m} = \sum_m d_m^* c_m$$
And, for the case of real solutions with negative Fourier coefficients equal to the complex conjugate of the postive Fourier coefficients, we have:
$$\braket{0|m} = d_0^* c_0 + 2 \cdot {\rm Re} \left( \sum_{m>0} d_m^* c_m \right) $$
Recall, that the python IRFFT library provides the Fourier coefficents with slightly different normalization:
\begin{equation}
A_m \equiv \frac{c_m}{\sqrt{a/N}} \exp\left(i \frac{2\pi m x_0}{a}\right)
\end{equation}
and
\begin{equation}
B_m \equiv \frac{d_m}{\sqrt{a/N}} \exp\left(i \frac{2\pi m x_0}{a}\right)
\end{equation}
so the dot product using the coefficients from IRFFT becomes:
$$\braket{0|m} = \Delta a \left\{ B_0^* A_0 + 2 \cdot {\rm Re} \left( \sum_{m>0} B_m^* A_m \right) \right\}$$
where, as before,
$$\Delta a = \frac{a}{N}$$.
In principle, we can substract any number of lower energy states to
find the next highest energy states, but round off errors will
eventually dominate.  Note that when subtracting multiple lower energy
states, you only need to normalize after subtracting the last
component.  One can combine the even versus odd trick to make this even more
efficient (e.g. subtracting the ground state from an even trial
function would lead to the 2nd excited state, if the energy
eigenstates alternate between even and odd.)

The numerical algorithm proceeds much the same way as in the previous
assignment, with one important modification: after {\bf each} random
variation of the wave function, you remove the lower energy
components, renormalize, and {\bf then} test for a decrease an energy.

\section{The Hydrogen Atom}

We can apply our numerical technique to find the radial wave functions of
the Hydrogen Atom.  In this case, we'll take our characteristic length scale $L$ to be the Bohr radius:
$$L \equiv \frac{4 \pi \epsilon_0 \hbar^2}{m e^2}$$
As always, this quantity should not be computed in your computer program, it will simply be set to one in every equation where it appears.
We continue to have our energy scale defined by:
$$\epsilon = \frac{\hbar^2}{2mL^2}$$
but with $L$ now defined we have an alternative form:
$$\epsilon = \frac{m}{2\hbar^2} \left( \frac{e^2}{4\pi\epsilon_0}\right)^2$$
Another useful quantity is:
$$\frac{e^2}{4\pi\epsilon_0} = 2 \epsilon L$$
Using this system of units, the protential for the Hydrogen atom becomes:
$$\frac{V(r)}{\epsilon} = - \frac{2L}{r}$$
The allowed energy levels are:
$$E_n =  \epsilon / n^2$$
where $n=1,2,3,\ldots$ is the principle quantum number.  The additional quantum numbers $l$ and $m$ refer to the angular momentum of the state.

\section{Radial Solutions}

Using the method of separation of variables (e.g. Griffiths Chapter 4) we write the total wave function as:
$$\Psi_{nlm}(r,\theta,\phi) = \frac{u_{nl}(r)}{r} Y_{lm}(\theta,\phi)$$
for quantum numbers $n,l,m$.  The radial part satisfies 
\begin{equation*}
  \frac{\hat{E}}{\epsilon} u(r) = -L^2\frac{d^2 u}{dr^2} + \left(-\frac{2L}{r}+ l(l+1) \frac{L^2}{r^2}\right)
\end{equation*}
Notice that this looks exactly like a 1-D particle, but with an effective potential:
$$\frac{V_{\rm eff}(r)}{\epsilon} = -\frac{2L}{r}+ l(l+1) \frac{L^2}{r^2}$$
The term proportional to $l(l+1)$ is called the centrifugal term, and
it accounts for the additional kinetic energy in the non-radial
direction, due to the angular momentum of the particle.

The first few radial solutions are:
$$u_{10} = 2 \cdot \frac{1
}{\sqrt{L}} \; \left( \frac{r}{L} \right) \; \exp\left(-\frac{r}{L}\right)$$
and:
$$u_{20} = \sqrt{2} \cdot \frac{1}{\sqrt{L}} \; \left( \frac{r}{2L} \right) \; \left(1 - \frac{r}{2L} \right)\; \exp\left(-\frac{r}{2L}\right)$$
and
$$u_{21} = \sqrt{\frac{2}{3}} \cdot {\sqrt{L}} \; \left( \frac{r}{2L} \right)^2 \; \exp\left(-\frac{r}{2L}\right)$$
Notice that none of these solutions are even or odd, but they do all vanish at $r=0$.

We can find these solutions using just a few modifications to our algorithm from the previous homework, via the replacement:
$$x \longrightarrow r$$
First, we must impose the boundary condition
$u(0)=0$
This is easy to implement in our Fourier expansion.  Because the sine vanishes at $x=0$, we will meet this condition provided:
$$c_0 = - 2 \cdot \sum_{m>0} {{\rm Re}\;{c_m}}$$
Recall also that for a real function $c_0$ is real.  So we only need to
modify our algorithm to not include $c_0$ in our random variations.
We fix the imaginary part of $c_0$ to zero, but after each variation,
we modify the real part of $c_0$ by the equation above.

In terms of the coefficients as provided by the IFFT, these only
differ by a constant real factor at $x=0$, so we have equivalently:
$$A_0 = - 2 \cdot \sum_{m>0} {{\rm Re}\;{A_m}}$$

We'll also need to change $x_0$ from $-a/2$ to $0$, appropriate for the radius $r$.


\section{Homework Problems}

\noindent
{\bf Please submit one PDF file (including output) and one python
  notebook file (.ipynb) of your completed assignment.}  If using
C/C++, submit text or PDF of your output and the source code.\\[3pt]


\noindent
{\bf Problem 1:} Building on your results from HW2, find the second
and third excited state of the Harmonic oscillator potential, and compare
the energy your determine for each with the expected value.\\[5pt]

\noindent
{\bf Problem 2:} For the hydrogen atom, using the known solutions,
plot $V(x)$, $|u_{10}(r)|^2$ and $|u_{20}(r)|^2$ in the same plot.
For better visiblity, offset the $y$-values by the energy of each
state.  So e.g. plot $u_{10}(x)=0$ at $V(x)/\epsilon=-1$.\\[5pt]

\noindent
{\bf Problem 3:} Find the Fourier coefficients of $u_{10}$ and
$u_{20}$, calculate their total energy from the Fourier coefficients,
and compare to the expected values.  Use this to validate the same
calculation you will be using in the variational method.\\[5pt]

\noindent
{\bf Problem 4:} Use the variational method to calculate the ground
state of the hydrogen atom.  Compare both the energy and the radial
wave function that you obtain to the known solution.  (I used $N=64$,
$A=20$ and a variation size of $0.1$ to achieve $10\%$ accuracy very
quickly.)\\[5pt]

\noindent
{\bf Problem 5:} Using the known solution, plot the effective radial potential
$V(x)$ for $l=1$ (including the centrifugal term!) and $|u_{21}(r)|^2$ in the same plot.\\[5pt]

\noindent
{\bf Problem 6:} Find the expectation value for the total energy for
$u_{20}$ and $u_{21}$ in the effective radial potential for $l=1$.
(Notice that this is not the correct effective potential for $u_{20}$
but calculate the expecation value anyway, it is instructive!)  You
should see that expectation value for $u_{20}$ in this potential is
higher than for $u_{21}$.  Do you understand why?  You should be able
to conclude now that $u_{21}(r)$ is the {\bf ground state} of the
effective radial potential for $l=1$.\\[5pt]

\noindent
{\bf Problem 7:} By varying the centrifugal term proportional, find
the energy for the excited states corresponding to $u_{21}$ and
$u_{32}$.  You {\bf do not} need to subtract out lower energy
components to do so!  Each of these excited states is in fact the
ground state of their effective potential.\\[5pt]

\noindent
{\bf Extra Challenge (No Credit):} If you had an easy time with the
above, see if you can also determine $u_{20}$.\\[5pt]


\end{document}
